% GENERATED FILE. Edit canonical lessons, then run npm run generate:books.
% canonical-source-hash: fnv1a64:dc4f426da5e4f2a3
% canonical-lessons: FR-C36-svo, FR-C36-et, FR-C36-ou-conj, FR-C36-mais, FR-C36-parce-que, FR-R36-deux-idees

\chapter{Two Ideas, One Sentence}
\label{ch:fr-joiners}
\hlchaptermodality{40}

\begin{chapteropening}
\textbf{By the end of this chapter:} \emph{I can name the order a French sentence holds to, and join two thoughts with and, or, but and because.}

Every verb the book owns is given a reason. The chapter also names the sentence shape the reader has been building for thirty-odd chapters without being told what it was, and says why French is stricter about word order than Latin ever was: Latin's endings said who did what, French wore them off, and position had to take the job.
\end{chapteropening}

\section[sujet, verbe, objet]{sujet, verbe, objet — the order that carries the meaning}

\label{lesson:FR-C36-svo}



\begin{quote}
You have been building French sentences for thirty-odd chapters and nobody has said what shape they are. Say \textbf{je parle français} and \textbf{parlez-vous français ?} and notice that only one of the two moves anything.

\begin{itemize}
\raggedright
  \item \emph{Recall:} two chapters back, the phrase that says a thing exists — \emph{il y a}
\end{itemize}
\end{quote}

\begin{cousinweb}
\begin{quote}
The doer comes first, the verb second, and what is done to comes third.
\end{quote}

\begin{quote}
\textbf{Marie mange le pain.} — Marie eats the bread.
\end{quote}

Turn those three around and you change who eats whom. French depends on the order, and the dependence has a cause.

Latin did not depend on it. \emph{Maria panem edit} and \emph{panem Maria edit} mean the same thing, because \textbf{panem} carries an ending that says \emph{this is the thing being eaten} wherever it stands. French wore those endings off — \emph{panem} became \emph{pain}, and the \emph{-em} that did the work went with it.

Once the endings were gone, something had to say who did what, and position took the job. So French is \textbf{stricter} about word order than its own parent was, and the strictness is the price of the simpler nouns.

Two places already break the order, and you own both: a question by inversion puts the verb first, and \emph{voici} and \emph{voilà} have no subject at all.
\end{cousinweb}

\subsection*{Your turn}
Say these aloud:

\begin{itemize}
\raggedright
  \item the plain order — \textquotedblleft{}Marie mange le pain.\textquotedblright{}
  \item turn it round and hear the meaning turn — \textquotedblleft{}Le pain mange Marie.\textquotedblright{}
  \item \textquotedblleft{}Je parle français.\textquotedblright{}
  \item the doer named by a possessive phrase — \textquotedblleft{}Le frère de Marie mange le pain.\textquotedblright{}
  \item the two shapes that break the order — \textquotedblleft{}Parlez-vous français ?\textquotedblright{} and \textquotedblleft{}Voici le café.\textquotedblright{}
\end{itemize}

\begin{tcolorbox}[breakable,colback=teal!4,colframe=teal!35!black,title={Before you move on}]
What order does a French sentence hold to? (\textbf{Doer, verb, done-to.}) Why is French stricter about it than Latin? (\textbf{Latin's endings said who did what; French wore them off}, so position had to take over.) Name the two shapes you own that break the order.
\end{tcolorbox}
\section[et]{et — \textquotedblleft{}and,\textquotedblright{} and the one \emph{t} that never wakes}

\label{lesson:FR-C36-et}



\begin{quote}
You have said \textbf{et vous ?} and \textbf{et toi ?} many times as glue in a conversation. It has never been a word of its own.
\end{quote}

\begin{cousinweb}
\begin{quote}
\textbf{et} — and.
\end{quote}

\begin{quote}
\textbf{du pain et du vin} — bread and wine.
\end{quote}

Latin \textbf{et}, unchanged in two thousand years, and the same word inside English \textbf{etcetera}.

It hides one genuine oddity. Every other final \textbf{t} in French comes back to life before a vowel — \emph{quand est-ce que} is \emph{kahn-TESS-kuh}, and you met that liaison with the numbers. The \textbf{t} of \textbf{et} never does:

\begin{quote}
\textbf{du pain et un café} — \emph{pan … ay … uh café}, with a clean break.
\end{quote}

Nobody can say why this one word refuses. It is old enough that grammarians record the exception and do not explain it, and it is worth knowing because the liaison you would naturally make there is the one liaison a French ear will notice as wrong.
\end{cousinweb}

\subsection*{Your turn}
Say these aloud:

\begin{itemize}
\raggedright
  \item \emph{et}, and the break after it — \textquotedblleft{}du pain et un café\textquotedblright{}
  \item the table joined — \textquotedblleft{}du pain et du vin\textquotedblright{}, \textquotedblleft{}du café et du lait\textquotedblright{}
  \item the household joined — \textquotedblleft{}mon frère et ma sœur\textquotedblright{}
  \item with the plural — \textquotedblleft{}des cafés et des amis\textquotedblright{}
\end{itemize}

\begin{tcolorbox}[breakable,colback=teal!4,colframe=teal!35!black,title={Before you move on}]
What Latin word is \emph{et}? (\emph{\textbf{Et}}, unchanged.) What does its \textbf{t} do before a vowel? (\textbf{Nothing} — it is the one final \emph{t} in French that never comes back.) Say \textquotedblleft{}bread and a coffee\textquotedblright{} with the break in the right place.
\end{tcolorbox}
\section[ou]{ou — \textquotedblleft{}or,\textquotedblright{} and the accent that is the only difference}

\label{lesson:FR-C36-ou-conj}



\begin{quote}
You own \textbf{où}, \textquotedblleft{}where.\textquotedblright{} Say it. This lesson takes the word that is spelled the same and means something else entirely.
\end{quote}

\begin{cousinweb}
\begin{quote}
\textbf{ou} — or. · \textbf{où} — where.
\end{quote}

\begin{quote}
\textbf{du café ou du thé ?} — coffee or tea?
\end{quote}

The two have nothing to do with each other. \textbf{ou} is Latin \textbf{aut}, \textquotedblleft{}or\textquotedblright{}; \textbf{où} is Latin \textbf{ubi}, \textquotedblleft{}where.\textquotedblright{} Two unrelated words wore down until they landed on the same three letters, and French put a grave accent on one of them to keep the page readable.

That is now the third such pair you have seen, and they work the same way every time: \textbf{a}/\textbf{à}, \textbf{la}/\textbf{là}, \textbf{ou}/\textbf{où}. The accent carries no sound at all in any of the three — it is punctuation for the eye.

Spoken French makes no distinction whatever. \emph{Où est le café ?} and \emph{ou est le café} are the same noise, and only the sentence settles it.
\end{cousinweb}

\subsection*{Your turn}
Say these aloud:

\begin{itemize}
\raggedright
  \item the choice — \textquotedblleft{}du café ou du thé ?\textquotedblright{}
  \item \textquotedblleft{}du pain ou du vin ?\textquotedblright{}, \textquotedblleft{}un chien ou un chat ?\textquotedblright{}
  \item the two joiners together — \textquotedblleft{}du pain et du vin\textquotedblright{}, \textquotedblleft{}du pain ou du vin\textquotedblright{}
  \item all three accent pairs — \textquotedblleft{}a/à, la/là, ou/où\textquotedblright{}
\end{itemize}

\begin{tcolorbox}[breakable,colback=teal!4,colframe=teal!35!black,title={Before you move on}]
Which Latin word is \emph{ou}? (\emph{\textbf{Aut}}.) And \emph{où}? (\emph{\textbf{Ubi}}.) Do they sound different? (\textbf{No} — the accent is for the eye.) Name the other two pairs the accent keeps apart. (\emph{\textbf{A/à}} and \emph{\textbf{la/là}}.)
\end{tcolorbox}
\section[mais]{mais — \textquotedblleft{}but,\textquotedblright{} and the Latin word for \emph{more}}

\label{lesson:FR-C36-mais}



\begin{quote}
\emph{Et} adds. \emph{Ou} offers a choice. Neither can say that the second half disagrees with the first.
\end{quote}

\begin{cousinweb}
\begin{quote}
\textbf{mais} — but.
\end{quote}

\begin{quote}
\textbf{Ça va bien, mais je suis fatigué.} — I'm well, but I'm tired.
\end{quote}

Said \emph{may}: the \textbf{ai} is the \emph{e} sound you met in \emph{lait}, and the \textbf{s} is silent.

\textbf{mais} is Latin \textbf{magis}, which meant \textbf{more} — and that is worth a moment, because the same word is still doing the older job next door. Spanish inherited \emph{magis} as \textbf{más} and it still means \emph{more}; French inherited it and drifted it into \emph{but}.

The road between the two is short. \emph{\textquotedblleft{}He is clever, more he is lazy\textquotedblright{}} is not far from \emph{\textquotedblleft{}he is clever, but he is lazy\textquotedblright{}} — a \emph{more} used to add a second, competing fact turns into a contrast if you use it that way for long enough. One Latin adverb, two daughters, two different parts of speech.
\end{cousinweb}

\subsection*{Your turn}
Say these aloud:

\begin{itemize}
\raggedright
  \item \emph{mais} — \emph{may}, the \emph{s} silent
  \item \textquotedblleft{}Ça va bien, mais je suis fatigué.\textquotedblright{}
  \item \textquotedblleft{}Il y a du café, mais il n'y a pas de thé.\textquotedblright{}
  \item the three joiners in a row — \textquotedblleft{}et, ou, mais\textquotedblright{}
\end{itemize}

\begin{tcolorbox}[breakable,colback=teal!4,colframe=teal!35!black,title={Before you move on}]
What did \emph{mais} mean in Latin? (\emph{\textbf{More}} — \emph{magis}.) Which language still uses that word for \emph{more}? (\textbf{Spanish}, as \emph{más}.) Say the three joiners.
\end{tcolorbox}
\section[parce que]{parce que — \textquotedblleft{}because,\textquotedblright{} and the three words inside it}

\label{lesson:FR-C36-parce-que}



\begin{quote}
You can ask \textbf{pourquoi ?}. Nothing you own can answer it.
\end{quote}

\begin{cousinweb}
\begin{quote}
\textbf{parce que} — because.
\end{quote}

\begin{quote}
\textbf{Pourquoi ? — Parce que je suis fatigué.}
\end{quote}

Three words are welded in there: \textbf{par} (\textquotedblleft{}by\textquotedblright{}), \textbf{ce} (\textquotedblleft{}this\textquotedblright{}) and \textbf{que} (\textquotedblleft{}that\textquotedblright{}). \emph{By this that…} — a phrase that pointed forward at the reason about to be given, said so often that the first two fused and the third stayed loose.

English built its own the same way. \textbf{Because} is \textbf{by} + \textbf{cause}, and the seam is still faintly visible if you look. Two languages, two prepositions, two pointing words, one job.

Before a vowel the \textbf{que} elides, exactly as \emph{le} and \emph{de} do:

\begin{quote}
\textbf{parce qu'il est tard} — because it's late.
\end{quote}

And the \textbf{ce} inside it is the same \emph{ce} you took two chapters ago — the fourth place that one small pointing word has turned up.
\end{cousinweb}

\subsection*{Your turn}
Say these aloud:

\begin{itemize}
\raggedright
  \item the question and its answer — \textquotedblleft{}Pourquoi ? — Parce que je suis fatigué.\textquotedblright{}
  \item elided — \textquotedblleft{}parce qu'il est tard\textquotedblright{}
  \item the three pieces, slowly — \emph{par … ce … que}
  \item all four joiners — \textquotedblleft{}et, ou, mais, parce que\textquotedblright{}
\end{itemize}

\begin{tcolorbox}[breakable,colback=teal!4,colframe=teal!35!black,title={Before you move on}]
What three words are inside \emph{parce que}? (\emph{\textbf{Par}}, \emph{\textbf{ce}}, \emph{\textbf{que}} — \textquotedblleft{}by this that\textquotedblright{}.) How did English build \emph{because}? (\textbf{By} plus \textbf{cause} — the same trick.) What happens to \emph{que} before a vowel? (\textbf{It elides} — \emph{parce qu'il}.)
\end{tcolorbox}
\section[Practice]{Deux idées — the chapter, run cold}

\label{lesson:FR-R36-deux-idees}



\begin{quote}
Say the four joiners. Then say what French lost that made its word order compulsory.
\end{quote}

\subsection*{Your turn: joined, chosen, contrasted, explained}
Say these aloud:

\begin{itemize}
\raggedright
  \item joined — \textquotedblleft{}du pain et du vin\textquotedblright{}, \textquotedblleft{}du café et du lait\textquotedblright{}
  \item chosen — \textquotedblleft{}du café ou du thé ?\textquotedblright{}, \textquotedblleft{}de l'eau ou du vin ?\textquotedblright{}
  \item contrasted — \textquotedblleft{}Ça va bien, mais je suis fatigué.\textquotedblright{}
  \item explained — \textquotedblleft{}Pourquoi ? — Parce que je suis fatigué.\textquotedblright{}
  \item all four in one breath — \textquotedblleft{}et, ou, mais, parce que\textquotedblright{}
\end{itemize}

\subsection*{Your turn: the distant band — every verb, given a reason}
Every verb this book owns can now carry a reason. Say the sentence, then say why.

Say these aloud:

\begin{itemize}
\raggedright
  \item \textquotedblleft{}Je parle français, parce que j'habite à Paris.\textquotedblright{}
  \item \textquotedblleft{}Je travaille, mais je suis fatigué.\textquotedblright{}
  \item the verbs of the page — \textquotedblleft{}Je lis et j'écris.\textquotedblright{}
  \item the verbs of the body — \textquotedblleft{}Je marche ou je cours.\textquotedblright{}
  \item \textquotedblleft{}J'entends, mais je ne comprends pas.\textquotedblright{}
  \item with the three modals — \textquotedblleft{}Je veux dormir, parce que je suis fatigué.\textquotedblright{} \textquotedblleft{}Je peux prendre le café.\textquotedblright{} \textquotedblleft{}Je dois travailler.\textquotedblright{}
  \item \textquotedblleft{}J'aime le café, mais je demande du thé.\textquotedblright{}
\end{itemize}

\begin{tcolorbox}[breakable,colback=teal!4,colframe=teal!35!black,title={Before you move on}]
Which joiner was Latin for \emph{more}? (\emph{\textbf{Mais}}, from \emph{magis}.) Which one has three words inside it? (\emph{\textbf{Parce que}}.) Which one's \emph{t} never makes a liaison? (\emph{\textbf{Et}}.) And which two are kept apart only by an accent? (\emph{\textbf{Ou}} and \emph{\textbf{où}}.)
\end{tcolorbox}
